% ============================================================================
% CAPÍTULO 2: MARCO TEÓRICO Y ANTECEDENTES
% ============================================================================
%
% REQUISITOS SEGÚN LA FACULTAD:
% 
% EL MARCO TEÓRICO DEBE INCLUIR:
% • Revisión de literatura para definir las variables (acorde a la pregunta de
%   investigación, relevante y actualizada)
% • Fuentes primarias
% • Un análisis crítico de la literatura y la información presentada corresponde
%   al producto individual de dicho análisis, estableciendo las conexiones pertinentes
%
% LOS ANTECEDENTES CONSIDERAN:
% • La historia del tema y/o problema de investigación en el contexto local y
%   nacional de forma cronológica (información presentada de forma organizada)
% • Investigaciones previas del problema de investigación (actuales)
% • Cuadro comparativo que integre los principales resultados de investigaciones
%   semejantes a la planteada, de acuerdo con la revisión de la literatura
%
% SISTEMA DE CITAS:
% • APA Séptima Edición: si tu tesis es de la LGAC en Formación Biomédica
% • Vancouver: si tu tesis es de la LGAC en Salud
% ============================================================================

\chapter{Marco Teórico y Antecedentes}
\label{cap:marco_teorico}

% ============================================================================
% PARTE 1: MARCO TEÓRICO
% ============================================================================

% ----------------------------------------------------------------------------
% SECCIÓN 2.1: Conceptos Fundamentales y Definiciones
% ----------------------------------------------------------------------------
\section{Conceptos Fundamentales y Definiciones}
\label{sec:conceptos_fundamentales}

% INSTRUCCIONES:
% - Define claramente las variables principales de tu investigación
% - Usa fuentes primarias (artículos científicos originales, no solo revisiones)
% - Incluye definiciones operacionales cuando sea necesario
% - Cita adecuadamente: \cite{clave} para APA o Vancouver según tu LGAC

[Define aquí las variables principales de tu investigación basándote en la literatura científica]

\subsection{[Variable 1]}
\label{subsec:variable1}

Según \cite{smith2023}, la \textit{[variable 1]} se define como \textit{[definición con fuente primaria]}. 

% EJEMPLO FUNCIONAL: Esta cita a smith2023 generará automáticamente la sección "Referencias"
% al final del documento. Puedes ver cómo se ve la referencia en formato APA 7. 

[Desarrolla el concepto, presenta diferentes perspectivas si las hay, y establece qué definición adoptarás para tu investigación]

\subsection{[Variable 2]}
\label{subsec:variable2}

[Repite la estructura para cada variable principal]

% ----------------------------------------------------------------------------
% SECCIÓN 2.2: Marco Teórico Principal
% ----------------------------------------------------------------------------
\section{Marco Teórico}
\label{sec:marco_teorico_principal}

% INSTRUCCIONES:
% - Presenta las teorías, modelos o enfoques que fundamentan tu investigación
% - Realiza un análisis crítico (no solo descripción)
% - Establece conexiones entre diferentes autores
% - Explica cómo estas teorías se relacionan con tu pregunta de investigación

\subsection{[Teoría o Modelo Principal]}
\label{subsec:teoria_principal}

La teoría de \textit{[nombre]} propuesta por \cite{bishop2006} establece que \textit{[descripción de la teoría]}.

% EJEMPLO: bishop2006 es un libro clásico de Machine Learning incluido en referencias.bib

[Análisis crítico: fortalezas, limitaciones, aplicaciones previas, relevancia para tu investigación]

\subsection{[Enfoques Complementarios]}
\label{subsec:enfoques_complementarios}

[Si hay otros enfoques teóricos relevantes, desarróllalos aquí]

% ----------------------------------------------------------------------------
% SECCIÓN 2.3: Análisis Crítico de la Literatura
% ----------------------------------------------------------------------------
\section{Análisis Crítico de la Literatura}
\label{sec:analisis_critico}

% INSTRUCCIONES:
% - No te limites a resumir estudios
% - Identifica patrones, contradicciones, vacíos
% - Establece conexiones entre diferentes trabajos
% - Posiciona tu investigación en el contexto de lo que ya se sabe

[Presenta tu análisis integrando diferentes fuentes]

Los estudios de \cite{garcia2022} y \cite{WHO2020} coinciden en que \textit{[punto de convergencia]}, sin embargo, difieren en \textit{[punto de divergencia]}, lo cual sugiere que \textit{[tu interpretación]}.

% EJEMPLOS FUNCIONALES: Estas referencias se formatearán automáticamente en APA 7

[Continúa construyendo tu argumento crítico]

% ============================================================================
% PARTE 2: ANTECEDENTES
% ============================================================================

% ----------------------------------------------------------------------------
% SECCIÓN 2.4: Historia del Tema (Contexto Local y Nacional)
% ----------------------------------------------------------------------------
\section{Antecedentes Históricos}
\label{sec:antecedentes_historicos}

% INSTRUCCIONES:
% - Presenta la evolución del tema de forma CRONOLÓGICA
% - Enfócate en contexto LOCAL (Chihuahua) y NACIONAL (México)
% - Organiza la información de manera clara
% - Conecta con el contexto internacional cuando sea relevante

\subsection{Contexto Internacional}
\label{subsec:contexto_internacional}

[Breve panorama internacional para establecer el estado general del conocimiento]

\subsection{Contexto Nacional (México)}
\label{subsec:contexto_nacional}

[Evolución del tema en México de forma cronológica]

\textbf{Década de [año-año]:}

En este período, \cite{autor_mexico} realizó \textit{[descripción]}...

\textbf{Década de [año-año]:}

[Continúa cronológicamente]

\subsection{Contexto Local (Chihuahua)}
\label{subsec:contexto_local}

[Si existe investigación local relevante, preséntala aquí]

% ----------------------------------------------------------------------------
% SECCIÓN 2.5: Investigaciones Previas
% ----------------------------------------------------------------------------
\section{Investigaciones Previas Relacionadas}
\label{sec:investigaciones_previas}

% INSTRUCCIONES:
% - Enfócate en investigaciones RECIENTES (últimos 5-10 años)
% - Prioriza estudios similares a tu pregunta de investigación
% - Identifica qué han hecho, qué han encontrado, qué limitaciones tienen
% - Prepara la información para el cuadro comparativo

\subsection{Estudios con Metodología Similar}
\label{subsec:metodologia_similar}

[Agrupa estudios por similitud metodológica o temática]

El estudio de \cite{autor2023} utilizó \textit{[metodología]} con una muestra de \textit{[n participantes]}, encontrando que \textit{[resultados principales]}. Sin embargo, este estudio presentó limitaciones en \textit{[limitaciones]}.

[Continúa con otros estudios relevantes]

% ----------------------------------------------------------------------------
% SECCIÓN 2.6: Cuadro Comparativo de Investigaciones
% ----------------------------------------------------------------------------
\section{Análisis Comparativo de Investigaciones}
\label{sec:cuadro_comparativo}

% INSTRUCCIONES:
% - Este cuadro es OBLIGATORIO según los requisitos de la facultad
% - Incluye: Autor/Año, Objetivo, Metodología, Muestra, Resultados, Limitaciones
% - Facilita identificar qué aporta tu investigación

La \Cref{tab:comparativa_investigaciones} presenta un análisis comparativo de las principales investigaciones relacionadas con el tema de estudio.

\begin{table}[htbp]
    \centering
    \caption{Cuadro comparativo de investigaciones previas}
    \label{tab:comparativa_investigaciones}
    \small
    \begin{tabular}{@{}p{2cm}p{2.5cm}p{2cm}p{1.5cm}p{3cm}p{2cm}@{}}
        \toprule
        \textbf{Autor (Año)} & \textbf{Objetivo} & \textbf{Metodología} & \textbf{Muestra} & \textbf{Resultados Principales} & \textbf{Limitaciones} \\
        \midrule
        
        Apellido et al. (2023) & 
        [Objetivo del estudio] & 
        [Tipo de estudio] & 
        n=[número] & 
        [Resultado 1, Resultado 2] & 
        [Limitación identificada] \\
        
        \midrule
        
        Apellido2 \& Apellido3 (2022) & 
        [Objetivo] & 
        [Metodología] & 
        n=[número] & 
        [Resultados] & 
        [Limitaciones] \\
        
        \midrule
        
        [Continúa con más estudios] &
        ... &
        ... &
        ... &
        ... &
        ... \\
        
        \midrule
        
        \textbf{Investigación Actual (Tu tesis)} & 
        [Tu objetivo general] & 
        [Tu metodología] & 
        n=[tu muestra] & 
        [Aportación esperada] & 
        [Limitaciones reconocidas] \\
        
        \bottomrule
    \end{tabular}
\end{table}

% ANÁLISIS DEL CUADRO:
\subsection{Síntesis del Análisis Comparativo}
\label{subsec:sintesis_comparativa}

Como se observa en la \Cref{tab:comparativa_investigaciones}, las investigaciones previas han abordado \textit{[aspectos]}. Sin embargo, ninguna ha \textit{[tu aportación diferencial]}, lo cual justifica la realización de la presente investigación.

Las principales brechas identificadas son:
\begin{enumerate}
    \item \textit{[Brecha 1]}
    \item \textit{[Brecha 2]}
    \item \textit{[Brecha 3]}
\end{enumerate}

% ============================================================================
% NOTAS FINALES PARA ESTE CAPÍTULO:
% ============================================================================
% 
% RECUERDA:
% 1. Todas las referencias deben ser VERIFICABLES
% 2. Prioriza FUENTES PRIMARIAS (artículos originales, no solo revisiones)
% 3. Realiza ANÁLISIS CRÍTICO, no solo descripción
% 4. El cuadro comparativo es OBLIGATORIO
% 5. Organiza la historia de forma CRONOLÓGICA
% 6. Cita en formato APA 7ma (LGAC Formación) o Vancouver (LGAC Salud)
%
% CITAS APA 7ma ejemplo:
%    Hernández, J. S., Tobón, S., & Vázquez, J. M. (2014). Estudio conceptual...
%
% CITAS Vancouver ejemplo:
%    1. Ferrero NA, Bortsov AV, Arora H, et al. Simulator training enhances...
%
% ============================================================================
